\documentclass[11pt,a4paper]{letter}

\usepackage[utf8]{inputenc}
\usepackage[T1]{fontenc}
\usepackage{mathpazo}
\usepackage[margin=1in]{geometry}
\usepackage{hyperref}
\usepackage{xcolor}

\hypersetup{
    colorlinks=true,
    linkcolor=black,
    urlcolor=blue!70!black
}

\signature{Alireza Moayedikia, Ph.D.\\
Swinburne Business School\\
Swinburne University of Technology\\
Melbourne, VIC, Australia\\
\href{mailto:amoayedikia@swinburne.edu.au}{amoayedikia@swinburne.edu.au}}

\address{Alireza Moayedikia\\
Swinburne Business School\\
Swinburne University of Technology\\
Melbourne, VIC, Australia\\
\href{mailto:amoayedikia@swinburne.edu.au}{amoayedikia@swinburne.edu.au}}

\longindentation=0pt
\date{\today}

\begin{document}

% \section*{Highlights}
\begin{itemize}
    \item ML-predicted failure risk parameterises Bernoulli service times for MRI scheduling.
    \item Prediction quality, not optimisation, drives scheduling performance gains.
    \item Risk-first heuristic matches MILP at 35{,}000$\times$ lower computation cost.
    \item Cognitive and behavioural features deliver 80\% of predictive signal over age.
    \item Age-based heuristics underperform schedules ignoring patient heterogeneity.
\end{itemize}


\end{document}

% \begin{letter}{%
% Editor-in-Chief\\
% \textit{Computers \& Operations Research}\\
% Elsevier}

% \opening{Dear Editor,}

% I am pleased to submit our manuscript entitled "Bridging Prediction and Prescription in Clinical Trial Scheduling" for consideration as an original research article in \textit{Computers \& Operations Research}. The paper develops a two-phase prescriptive analytics framework that integrates machine learning prediction with stochastic optimization for appointment scheduling in capacity-constrained diagnostic services. Grounded in Alzheimer's disease clinical trials, we train gradient boosting models on cognitive assessments, behavioral symptoms, and clinical staging to estimate patient-specific MRI scan failure probabilities, then embed these predictions as Bernoulli service time distributions in a stochastic job-shop scheduling formulation.

% \medskip

% Computational experiments on simulation-augmented clinical cohorts yield three results with broad implications for the predict-then-optimize paradigm:
% \begin{enumerate}
%     \item \textit{Prediction dominates optimization.} A simple longest-processing-time heuristic achieves statistically indistinguishable performance from a mixed-integer program, despite $35{,}000\times$ greater computation time for the latter. This challenges the conventional emphasis on optimization sophistication.
%     \item \textit{Feature richness dominates model complexity.} Machine learning applied to demographic features alone provides no benefit over simple age-bracket rules. Value emerges only when cognitive and behavioral biomarkers are incorporated---approximately 80\% of predictive signal resides in non-demographic features.
%     \item \textit{Partial information can be worse than no information.} Age-based heuristics, which embody common clinical domain knowledge, actually underperform deterministic scheduling that ignores patient heterogeneity entirely.
% \end{enumerate}

% \medskip
% The manuscript directly addresses several themes central to the journal's scope. It contributes to healthcare scheduling optimization, a foundational application of operations research, by introducing patient-specific stochastic service time distributions parameterized through machine learning---bridging two domains the journal has actively promoted. The paper engages with recent advances in data science and machine learning within OR, including the predict-then-optimize paradigm surveyed in recent EJOR publications. Methodologically, the paper combines stochastic mixed-integer programming with gradient boosting classification, and provides rigorous computational experiments with seven baseline comparisons across 100 replications---ensuring the numerical experiments carry scientific value in their own right. The finding that prediction quality, not optimization sophistication, drives scheduling performance offers a practically important insight for the OR community.

% \medskip
% \textbf{Declarations.}\quad
% This manuscript has not been published previously and is not under consideration elsewhere. All data sources are acknowledged. The author declares no competing financial interests. The use of AI-assisted tools in manuscript preparation is disclosed in the manuscript per Elsevier's policy.

% \medskip
% Thank you for considering this submission. I look forward to your response.

% \closing{Sincerely,}

% \end{letter}
